\section{Time: serial computation and the width of its carrier}
\label{sec:time}

The complaint addressed here is that chain-of-thought is a workaround for a
model that never learned to reason latently. We show the strong form is false
under standard assumptions (Section~\ref{sec:time-necessary}); that the
comparison between media is a genuine trade-off with a computable crossover, but
in the \emph{opposite} direction to the one usually asserted
(Section~\ref{sec:time-capacity}); and that an unconditional separation is
available only in a restricted read model, with a barrier explaining why
(Section~\ref{sec:time-sep}).

\subsection{Serial computation is load-bearing}
\label{sec:time-necessary}

\begin{theorem}[Depth is not a substitute for steps]
\label{thm:tc0}
Let $\mathcal{T}$ be the class of languages decided by fixed-depth,
log-precision transformer decoders whose feedforward maps are computable in
space linear in their input, and let $\mathcal{T}_{\mathrm{CoT}}$ be the class
decided by the same decoders granted $\mathrm{poly}(n)$ chain-of-thought steps.
Then $\mathcal{T} \subseteq$ uniform $\mathrm{TC}^0$ \cite{merrill2023} and
$\mathcal{T}_{\mathrm{CoT}} = \mathrm{P}$ \cite{merrill2024}. Consequently, if
$\mathrm{TC}^0 \ne \mathrm{P}$, then
$\mathcal{T} \subsetneq \mathcal{T}_{\mathrm{CoT}}$.
\end{theorem}

\begin{remark}[Three hypotheses that are usually dropped]
\label{rem:tc0-care}
First, the uniformity in \cite{merrill2023} is a condition on the \emph{model}
--- the feedforward maps must be space-constructible --- not a free adjective on
$\mathrm{TC}^0$; without it the containment is in non-uniform $\mathrm{TC}^0$ and
the separation required becomes $\mathrm{TC}^0/\mathrm{poly} \ne \mathrm{P}$.
Second, these are statements about \emph{classes of languages}, not about a
particular network: a single trained model decides one language, so
``this model cannot do it without CoT'' does not follow from
Theorem~\ref{thm:tc0} for any individual model, only the non-existence of
\emph{some} fixed-depth decoder for the hard language. Third, the ``$=\mathrm{P}$''
direction additionally requires the generalized pre-norm architecture of
\cite{merrill2024}. With these in place the conclusion still stands, and it is
fatal to the strong form of the complaint: serial steps buy computational power
that no amount of latent skill at fixed depth recovers.
\end{remark}

\subsection{Scratchpad capacity: a trade-off, not a dominance}
\label{sec:time-capacity}

We now compare media. The natural intuition --- a token carries
$\log_2|\Vocab| \approx 17$ bits, a hidden state carries thousands, therefore
latent recurrence transports far more per step --- is correct about
\emph{per-step width} and wrong about \emph{total state}, because the two media
differ in whether that width accumulates.

\begin{definition}[Step medium]
\label{def:medium}
A \emph{step medium} of width $b$ is a mechanism by which one serial step
transmits information to subsequent steps. It is \emph{accumulating} if the
contents written at every step remain readable at all later steps, and
\emph{non-accumulating} if each step's write replaces the previous contents.
\end{definition}

Chain-of-thought is accumulating with $b = \log_2 V$: the emitted tokens persist
in the context and every later step attends over all of them. A looped or
depth-recurrent model carrying a single state $s_i \in \R^{d}$ is
non-accumulating with $b = d\,b_{\mathrm{eff}}$, where $b_{\mathrm{eff}}$ is the
effective bits per coordinate. A latent model that \emph{appends} each
continuous thought to the sequence, as in \cite{hao2024coconut}, is accumulating
with $b = d\,b_{\mathrm{eff}}$.

\begin{theorem}[Scratchpad capacity]
\label{thm:capacity}
After $n$ steps, the set of distinguishable carrier contents has cardinality at
most $2^{nb}$ for an accumulating medium of width $b$, and at most $2^{b}$ for a
non-accumulating one.
\end{theorem}

\begin{proof}
An accumulating medium's contents after $n$ steps are the tuple of the $n$
writes, each drawn from a set of size at most $2^{b}$. A non-accumulating
medium's contents are the last write alone, from a set of size at most $2^b$;
iterating a map cannot increase cardinality, since $|G(A)| \le |A|$ for any $G$
and any $A$.
\end{proof}

\begin{corollary}[Crossover]
\label{cor:crossover}
Chain-of-thought carries strictly more state than a fixed-size latent loop once
\begin{equation}
  n \;>\; n^{*} \;=\; \frac{d\,b_{\mathrm{eff}}}{\log_2 V} .
\end{equation}
For $d = 1024$, $b_{\mathrm{eff}} = 4$ and $V = 151936$ this is
$n^{*} \approx 238$ steps. Against an \emph{accumulating} latent medium of the
same width, by contrast, chain-of-thought is dominated at every $n$.
\end{corollary}

\begin{remark}[This reverses the usual claim, and removes a contradiction]
\label{rem:reversal}
The informal argument --- $17$ bits against thousands, therefore latent
recurrence wins --- silently compares a per-step width with a per-step width
while the quantity that matters for a long derivation is total accessible state.
A fixed-size latent loop \emph{overwrites}; a chain of tokens \emph{appends}.
Past a few hundred steps, which is exactly the regime long chains of thought
occupy, the token stream is the larger memory. This is the same phenomenon as
Theorem~\ref{thm:recall} viewed from the other side: bounded state is a
liability for long-horizon work, and it does not stop being one because the state
is continuous. The defensible form of the complaint is therefore narrower and
sharper than the original: \emph{the discretisation is a real cost per step, and
it is worth paying only where the medium accumulates}. Latent recurrence should
append, not overwrite --- and a model that appends continuous thoughts dominates
chain-of-thought on this measure at every horizon.
\end{remark}

\subsection{Simulation, and where a separation can live}
\label{sec:time-sep}

\begin{theorem}[Latent recurrence simulates chain-of-thought]
\label{thm:sim}
Fix a model $M$ with $L$-Lipschitz blocks, a finite prompt set, and a decoding
temperature $\tau > 0$, and suppose the logit margin at every step of the
trajectory is at least $\delta > 0$. Then there is an $n$-step latent-recurrent
computation (Definition~\ref{def:latent}) whose induced token sequence equals
the $n$-step chain-of-thought trajectory of $M$ exactly.
\end{theorem}

\begin{proof}[Proof sketch]
Take $G$ to be: run $M$'s stack on the current state, apply $W_U N(\cdot)$, form
$\softmax(\cdot/\tau)$, and write back the corresponding convex combination of
embeddings. Under a uniform margin $\delta$, taking $\tau$ small enough makes
this softmax within any desired $\varepsilon$ of the one-hot at the argmax, and
finiteness of the prompt set makes the required $\tau$ uniform; the Lipschitz
hypothesis converts the $\varepsilon$ on the written state into a bound on the
downstream deviation, which stays below $\delta/2$ and hence preserves every
subsequent argmax. Each $G$ is a composition of $M$'s own maps, so the
computation lies in the latent-recurrent class.
\end{proof}

\begin{remark}[A precision tension]
Theorem~\ref{thm:sim} requires $\Omega(n\log L)$ bits of precision to keep the
margin through $n$ steps, whereas Theorem~\ref{thm:tc0} is stated for
log-precision models. The two results are therefore about different machines,
and no single model satisfies both hypotheses for large $n$. We state this
rather than paper over it.
\end{remark}

\begin{remark}[Strictness does not follow]
\label{rem:nostrict}
Theorem~\ref{thm:sim} gives an inclusion. The converse separation does
\emph{not} follow from Theorem~\ref{thm:capacity}, which after
Corollary~\ref{cor:crossover} runs the other way for large $n$; and a naive
counting separation fails outright, because intermediate states of a
deterministic task are functions of the prompt. Concretely, a prompt containing
two permutations of $S_{30}$ induces $30!$ distinct ``intermediate
configurations'', which a counting argument would say requires seven CoT tokens
to distinguish, while a model in fact distinguishes them at $n=0$ by attending to
the prompt.
\end{remark}

The escape in Remark~\ref{rem:nostrict} is re-reading the prompt, and it must be
closed by hypothesis rather than by choice of task.

\begin{theorem}[Separation under a streaming read schedule]
\label{thm:stream}
Let $\mathrm{PROD}_{N,T}$ be the problem of composing $T$ given elements of
$S_N$. Consider a \emph{step machine} that reads one input symbol per step, may
perform arbitrary computation within a step, and carries information between
steps only through an interface of width $b$ bits. Such a machine solves
$\mathrm{PROD}_{N,T}$ for all $T$ if and only if $b \ge \log_2 (N!)$.
\end{theorem}

\begin{proof}
Sufficiency: carry the running product, which is an element of $S_N$ and so fits
in $\lceil \log_2 N!\rceil$ bits; each step composes it with the incoming
generator. Necessity: the induced map from input prefixes to interface contents
must be constant on Myhill--Nerode classes of the language
$\{\,(g_{1:T}, \pi) : g_T\cdots g_1 = \pi\,\}$. Two prefixes with distinct
running products are inequivalent, since a suitable suffix distinguishes them,
so the machine needs at least $|S_N| = N!$ interface states.
\end{proof}

\begin{corollary}
\label{cor:instantiate}
With $V = 151936$ and $d\,b_{\mathrm{eff}} = 4096$, a one-token-per-step CoT
interface fails from $N = 9$ onwards, since
$\log_2 8! = 15.30 < \log_2 V = 17.21 < 18.47 = \log_2 9!$, while the latent
interface suffices to $N = 536$.
\end{corollary}

\begin{proposition}[The streaming hypothesis is doing real work]
\label{prop:confound}
With full access to the prompt and unbounded computation within a step, a
\emph{single}-step machine solves $\mathrm{PROD}_{N,T}$ for every $T$.
\end{proposition}

\begin{proof}
The entire input is available in one step; compute the product.
\end{proof}

\begin{remark}[A barrier]
\label{rem:barrington}
Proposition~\ref{prop:confound} means any separation in the full-attention
setting must be a circuit-complexity statement, not a counting one --- and such
statements are currently out of reach. By Barrington's theorem, word problems
over $S_5$ are $\mathrm{NC}^1$-complete, so a theorem of the form
``$\mathrm{PROD}_{N,T}$ requires at least two chain-of-thought steps for a
fixed-depth transformer'' would imply $\mathrm{TC}^0 \ne \mathrm{NC}^1$. The
bandwidth separation of Theorem~\ref{thm:stream} is therefore the strongest
unconditional result available, and it is conditional on a read schedule that
real architectures do not obey. Our experimental protocol
(Section~\ref{sec:unified}) must enforce that schedule rather than assume it.
\end{remark}

\subsection{Discretisation as error correction}

\begin{openproblem}[Bandwidth against drift]
\label{op:drift}
Fix a latent update $G$ and per-step perturbation of variance $\eta^2$, and let
$\mathcal{C}$ be a codebook of rate $R$ bits per step onto which the state is
projected every $k$ steps. Quantify the largest $n$ for which the trajectory
remains recoverable, as a function of $(\eta, R, k)$.
\end{openproblem}

The premise that discretisation corrects errors is not free: projection reduces
error only under a margin condition, and near a cell boundary quantisation
\emph{amplifies} perturbation. For the tractable case of an isometric $G$,
Gaussian noise of variance $\eta^2$ per step, and a lattice codebook of rate $R$
in dimension $d$, standard lattice-quantisation arguments give a crossover
depth scaling as $(P/\eta^{2})\,2^{-2R/d}$ for signal power $P$, below which the
unquantised trajectory is recoverable and above which projection wins. We state
the general case as open because neither the isometry nor the noise model
survives contact with a trained network. The design question it implies --- how
often to project, onto a \emph{learned} codebook rather than the vocabulary --- is
the one parameter in this paper that no published architecture sweeps.
